% ---------------------------------------------------------------------------
% The corpus, the two grammars and the translator, compressed to what the measurement
% needs. The instantiation paragraph is load-bearing in this version rather than
% incidental: it is the first place the paper's result shows up, as a parse rate.
% Tables are the complete paper's.
% ---------------------------------------------------------------------------

\section{The Corpus, the Grammars, and the Translator}
\label{s2:sec:translator}

\subsection{The corpus}
\label{s2:sec:corpus}

Three sources, each pinned to a commit or a retrieval date, chosen to put a vendor's own
queries and a community's queries on the same footing. $41$ compiled
\texttt{kube-prometheus} dashboards supply PromQL written and maintained by the Kubernetes
monitoring projects themselves, pinned to its rendered form~\citep{kubeprom,kubemixin};
the $199$ most-downloaded Grafana community dashboards supply PromQL written by everyone
else~\citep{grafanadash}; and the $283$ integration dashboards of Datadog's
\texttt{integrations-core} supply Datadog queries written by Datadog, which matters
because no reference grammar for that language is published and vendor-authored queries
are the closest thing to a specification by example. Extraction yields $22{,}903$ query
strings; nine turn out to be LogQL --- Loki panels in dashboards registered under a
Prometheus datasource --- and are dropped rather than counted as PromQL the grammar fails
on, leaving $22{,}894$ in three populations: $8{,}770$ Datadog metric queries, $8{,}102$
PromQL expressions and $6{,}022$ Datadog formula expressions, the last held aside.
Everything below is over the whole of it, with nothing sampled and nothing
hand-classified.

\paragraph{Dashboard queries are not queries.}
A panel ships \texttt{\$pod}, \texttt{\$\_\_rate\_interval} or \texttt{\$host} where a
value belongs, so the first measurement this paper made was of its own instrument failing:
$11.8\%$ of distinct Datadog and $71.2\%$ of distinct PromQL queries parse as shipped,
against $99.9\%$ and $99.1\%$ once instantiated.
\texttt{corpus.normalise} performs the substitution a dashboard performs --- a PromQL range
variable becomes a duration, a Datadog scope variable becomes \texttt{*}, a variable
anywhere else outside a string literal becomes an identifier, and a variable \emph{inside}
one is left alone, because \texttt{\{pod=\~{}"\$pod"\}} already parses --- and every
measurement in this paper is made on the instantiated form.

That is a methodological note and also the paper's finding arriving nine sections early.
The object a dashboard migration translates is a query \emph{schema}, and a translator that
takes query strings is solving a different problem from the one its users have. The
substitution above is a stand-in for a dashboard's own resolution, performed by us because
the dashboards are not running; \S\ref{s2:sec:unsaid} is what happens when the same
observation is made about values no substitution can supply.

\subsection{The grammars exist, and they are deterministic}

Neither vendor publishes a grammar. We infer one for each language, and the Datadog
grammar is, as far as we have found, the first published for it --- which is a statement
about a search and is reported as one. Both grammars admit deterministic one-token-lookahead
parsers, at the sizes Table~\ref{s2:tab:grammars} reports. Datadog's is conflict-free.
PromQL's carries two shift/reduce conflicts on a left parenthesis, both resolved the way
the reference server resolves them; the complete paper documents the resolution and shows
it is the one the published implementation makes, and declines to claim that no LALR(1)
grammar for PromQL exists. That matters here only because it licenses what follows: a
unique parse tree per query is what makes a syntax-directed translation well defined at
all, and a table-driven parse is what makes a corpus of this size processable rather than
sampleable.

\begin{table}[htbp]\centering\footnotesize
\begin{tabular}{@{}llrrrl@{}}
\toprule
Language & Grammar & Terminals & Rules & LALR states & Conflicts \\
\midrule
promql & v1 & 32 & 73 & 130 & 2$\times$ \texttt{Shift/Reduce on LPAR} \\
promql & v2 & 33 & 82 & 145 & 2$\times$ \texttt{Shift/Reduce on LPAR} \\
datadog & v1 & 19 & 35 & 66 & none \\
datadog & v2 & 22 & 60 & 111 & none \\
\bottomrule
\end{tabular}

\caption{The two inferred grammars, at the two revisions this work produced for each.
Conflicts are LALR(1) conflicts reported by the generator; Datadog's grammar has none at
either revision, and PromQL's two are the shift/reduce pair on a left parenthesis that the
reference server resolves the same way.}
\label{s2:tab:grammars}
\end{table}

\noindent
Those instantiated rates are the size of a fit rather than a result, and
\S\ref{s2:sec:threats} reads them as one: grammar v2 is v1 plus every production the corpus
showed was missing, so the climb from $92.9\%$ to $99.9\%$ on the Datadog side is what
reading the corpus bought. The raw rate, on the other hand, is not a property of either
grammar at all. A Datadog string parses as shipped exactly when it carries no template
variable, with two exceptions in the whole held-out set, so $11.8\%$ is the
template-variable share counted the other way round.

\subsection{The translator}

$T$ factors as $\pi_1\circ E\circ L$ through an intermediate representation, as
Definition~\ref{s2:def:translator} sets out. The lowering reads a Datadog query into the
IR; the emission writes PromQL and returns, beside the expression, the set of caveats
recording what it had to assume. Every refusal is a $\bot$ with a gap code. The IR is the
only object in this work whose denotation is written out, and the exactness results of the
complete paper are proved over $E(L(q))$ rather than over $T(q)$, because they need the
caveat set that $E$ returns alongside the expression.

One of the complete paper's results is worth naming, because \S\ref{s2:sec:coverage} rests
on it rather than on inspection: a kernel soundness theorem covering $95.7\%$ of the
corpus, with six hypotheses and two side conditions, three of which the emitter did not
record when the theorem was first stated and now returns as codes. An unrecorded
precondition is free in every number a paper reports, which is exactly why they accumulate,
and each of those three turned into a row of \S\ref{s2:sec:unsaid}'s table without moving
$\kappa$ by a single occurrence.
